\section{Flavor boundary and representation algebra}
\label{sec:flavour}

The electroweak mass-ratio problem and the flavor problem carry separate obligations in this manuscript.  The working boundary condition is:
\begin{equation}
\begin{aligned}
  &\text{DeVries electroweak determinant}\\
  &\qquad\text{and generation/flavor organization}\\
  &\text{are linked only by an explicit model}.
\end{aligned}
  \label{eq:flavor-boundary}
\end{equation}
This prevents a compactification or endpoint picture from turning a useful flavor pattern into an electroweak derivation by naming alone.

Representation theory remains essential.  In the Standard Model and in the familiar grand-unification lineage, fermions live in representations of the internal gauge group, while gauge bosons live in the adjoint sector \cite{BaezHuerta2010,Britto2021}.  SU(5), Spin(10), Pati--Salam, and \(E_6\) show how a larger algebra can unify the Standard Model representation data.  This is the correct role of the GUT material in the present paper: it supplies representation discipline and a comparison class for structural weak-angle statements.

The DeVries construction uses a different representation question.  The electroweak determinant asks why one sample is associated with the Higgs/order-parameter doublet invariant and the other with the adjoint/current invariant:
\begin{equation}
  J_H=T_H(T_H+1)=3/4,\qquad
  J_{\rm adj}=T_{\rm adj}(T_{\rm adj}+1)=2.
  \label{eq:flavor-j-boundary}
\end{equation}
Flavor representations can constrain a completion.  The missing step is a derived map from endpoint, compactification, or localized matter data into the two-channel electroweak operator.

The SO(32)-flavor idea should therefore be treated as an organizing boundary sector.  In an endpoint language it may label external or internal flavor degrees of freedom, Chan--Paton-like data, or a bookkeeping algebra for fermion multiplicities.  Its required checks are concrete:
\begin{enumerate}
\item identify the representation space carrying the SO(32)-flavor labels;
\item map the Standard Model chiral fermions into that space;
\item show anomaly cancellation and hypercharge normalization;
\item state the relation to the global quotient \(\Gamma\subset\mathbb Z_6\);
\item couple the flavor sector to the electroweak determinant only through a derived operator.
\end{enumerate}

This keeps the flavor idea useful while protecting the central claim.  The electroweak ratio concerns W/Z pole structure and the Higgs/order-parameter channel.  Generation structure concerns chiral matter, Yukawa textures, and possible endpoint or compactification multiplicities.  A string/Kaluza--Klein completion may eventually bind them together.  The binding has to be an equation or a local geometric construction.

The compactification version has an additional caveat.  Geometry can generate gauge symmetries, chiral localization, and multiple wavefunctions with the same gauge quantum numbers \cite{WittenKK1981,AcharyaWitten2001,BraunEtAl2019}.  That makes topology relevant to flavor.  The manuscript still treats generation counting as a separate task because the DeVries determinant is a two-channel electroweak object.  The decisive test is whether the same local model that yields generations also yields
\begin{equation}
  \det(\lambda-\mathcal O_J)=\lambda^2+J\lambda-J
  \label{eq:flavor-det-test}
\end{equation}
with the ordered electroweak samples in Eq.~\eqref{eq:flavor-j-boundary}.

The practical rule for later sections is therefore simple.  Flavor material may enter as a constraint on endpoints, compactification, anomaly cancellation, and global form.  It may motivate a local operator.  The W/Z mass-ratio argument rests on the electroweak determinant and its pole placement.

\paragraph*{Section status.}  This section states a source-supported boundary condition.  The SO(32)-flavor interpretation remains conjectural until a representation map and coupling to \(\mathcal O_J\) are supplied.
